% 正赛文稿模板：每方一份。下面的十节是中科大新生赛的稿件清单，只是示例；写之前跑
%   python3 scripts/format_info.py <赛制> --headings 正方   （反方同理）
% 用它输出的标题原样替换，多的节删掉、少的节按 references/stage-playbooks.md 同类型的写法补。
% 括号里的字数用 scripts/check_speeches.py 数出来再填。质询、盘问、对辩、自由辩是 stage，不计字数。
% 稿件正文里：重音用 \stress{…}（每篇最多两处），舞台提示用 \aside{…}，临场位用 \reserve{…}，都不计字数。
\documentclass[script,pro]{debate}   % 正方 pro，反方 con
\motion{<辩题>}
\stance{正方}{<持方表述>}
\meta{赛制}{<…>}
\meta{口径来源}{备赛文档 第十节 <正方 / 反方> 口径表}
\begin{document}
\debatetitle
\begin{thesisbox}[全场一句话立论]<…>\end{thesisbox}

\begin{speech}{1. 正一开篇立论稿}{正文 NNN 字 / 180 秒}
谢谢主席。<定义 → 中立判准与双方义务 → 论点一（主张 → 机制 → 论据 → 回扣）→ 论点二 → 抛出后场问题 → 价值一句 → 收束>
\end{speech}
\begin{contingency}
\ifthen{<对方定义与我方不同>}{<30 秒定义之争口径>}
\ifthen{<时间不够>}{<删哪一句>}
\end{contingency}

\begin{stage}{2. 正一被质询防守预案}{反四质询，120 秒双边计时}
\textbf{原则}：<回答短、守住前提、不否认常识>
\begin{qa}
\qarow{<对方会问>}{<我方回答（简短，守住前提）>}
\end{qa}
\end{stage}

\begin{stage}{3. 正四质询反一问题链}{120 秒，双边计时}
\textbf{总目标}：<拿到什么承认 / 暴露什么矛盾>
\begin{chain}{链一：<名称>}{<目标>}
\q{<问题，25 字内>}{答“是”：<下一问>；答“否”：<下一问>}
\q{<…>}{<…>}
\cut{<打断用语>}
\close{<收束语>}
\end{chain}
\begin{chain}{链二：<名称>}{<目标>}
\q{<…>}{<…>}
\end{chain}
\end{stage}

\begin{speech}{4. 正二驳论稿}{正文 NNN 字 + 临场位 NNN 字 / 120 秒}
<总起 → 驳点一 → 驳点二 → 收束；反方先驳，正方后驳并预留约 100 字临场位>
\reserve{<临场回应反二的位置与口径，约 100 字>}
\end{speech}
\begin{contingency}
\ifthen{<…>}{<…>}
\end{contingency}

\begin{stage}{5. 正二对辩要点}{90 秒，单边计时}
\begin{points}{进攻点（4–5 个，每个 15–20 秒）}
\item <…>
\end{points}
\begin{points}{备用点（6–8 个）}
\item <…>
\end{points}
\begin{points}{对方最可能问的三个问题 → 回应}
\item <…>
\end{points}
\end{stage}

\begin{stage}{6. 正三盘问问题链}{90 秒，单边计时}
\begin{chain}{链一：<名称>}{问二辩 → 问四辩同一问题，检验口径}
\q{（问反二）<…>}{<…>}
\q{（问反四）同样的问题}{<两人若不一致，小结重锤>}
\end{chain}
\begin{chain}{链二：<名称>}{<目标>}
\q{<…>}{<…>}
\end{chain}
\textbf{盘问目标}：<小结里要证明什么>
\end{stage}

\begin{stage}{7. 正二、正四被盘问防守预案}{两人口径必须一致}
\begin{qaa}
\qaarow{<对方会问>}{<二辩答>}{<四辩答（必须一致）>}
\end{qaa}
\end{stage}

\begin{speech}{8. 正三盘问小结稿}{正文 NNN 字 / 90 秒}
<盘问成果 → 与论点的连接 → 当前战场总结；反方小结预留约 80 字临场位并在标题里写「+ 临场位 NNN 字」>
\end{speech}

\begin{stage}{9. 自由辩战场设计}{240 秒}
\begin{battleground}{战场一：<名称>（前 90 秒，主战场）}
\begin{fields}
\field{目标}{<…>}
\field{开场问题}{<封闭式，20 字内>}
\field{对方答 A → 追问}{<…>}
\field{对方答 B → 追问}{<…>}
\field{对方可能反攻 → 我方防守}{<…>}
\field{收束语}{<…>}
\end{fields}
\end{battleground}
\begin{battleground}{战场二：<名称>（中 90 秒）}
<同上>
\end{battleground}
\begin{points}{必问（前两分钟内问完）}
\item <…>
\end{points}
\begin{points}{必答（15 秒以内正面回答）}
\item <问题 → 答案>
\end{points}
\textbf{时间分配与站位}：<…>
\end{stage}

\begin{speech}{10. 正四结辩稿}{正文 NNN 字 / 210 秒，留 30–40 秒临场回应}
<归纳分歧 → 处理对方最强点 → 我方论点收束（回到一辩稿抛出的问题）→ 价值升华（具体的人与时刻，不是宏大叙事）→ 结尾短句。反四写封路式，标题用「封路式」；正四留 30–40 秒临场位>
\end{speech}
\begin{contingency}
\ifthen{<…>}{<…>}
\end{contingency}

% 有奇袭环节的赛制（如 2025 校赛）才有第 11 节；其他赛制整节删掉。
% \begin{stage}{11. 正方奇袭预案}{质询 150 秒双边 / 申论 120 秒，二选一}
% \textbf{发动判断}：<什么条件下发动、在哪个环节之后、由谁喊、暗号是什么>
% \begin{chain}{奇袭质询链一：<通用链，打定义或判准>}{<目标>}
% \q{<…>}{<…>}
% \end{chain}
% \begin{speech}{奇袭申论稿}{正文 NNN 字 / 120 秒}
% <开头留一句接场上刚发生的事，正文重建被打塌的论点或立起未被触碰的战场>
% \end{speech}
% \textbf{不发动的情况}：<…>
% \end{stage}
\end{document}
